
In the following results are for a regar regression model specified in the following way:

We consider a set of $N$ observations generated from a spatial regression
framework.  Let $y$ denote the response vector, $X = (X_1, X_2, X_3)$ the
explanatory variables, and $W$ a spatial weights matrix (Rook, Queen, or
$k$-nearest neighbors with $k=6$).  The spatial Durbin model (SDM) used in the
simulation study is given by
\begin{equation}
	y = \rho Wy + X\beta + \varepsilon,
	\label{eq:sdm}
\end{equation}
where $\rho$ is the spatial autoregressive parameter, $\beta$ are the direct effects

The error term follows
\[
	\varepsilon \sim \mathcal{N}(0, \sigma^2 I_N).
\]

The true parameter values in the data-generating process are
\[
	\beta = (4,5,6,7)^{\top}, \qquad \rho = 0.7.
\]

Looking over the results we see that the best preforming model is the Queens
model as this is the model used to generate the data. The second best preforming
model is the penalized tensor products. This shows that if there is no reasonable
argument for picking the spatial weights matrix the penalized tensor products are a
reasonable starting point for controlling for spatial variability.

\begin{table}[H]
	\centering
	\caption{Simulation results for 1000 runs comparing spatial specifications}
	\begin{tabular}{lcccccc}
		\hline
		 & \textbf{Rook} & \textbf{Queen} & \textbf{KNN (6)} & \textbf{Tensor} & \textbf{Base} \\
		\hline
		Intercept
		 & 3973.28
		 & 0.40
		 & 92.49
		 & 8252.62
		 & 7177.34                                                                             \\

		$X_1$
		 & 0.74
		 & 0.0032
		 & 0.30
		 & 0.21
		 & 1.22                                                                                \\

		$X_2$
		 & 0.90
		 & 0.0032
		 & 0.30
		 & 0.26
		 & 1.57                                                                                \\

		$X_3$
		 & 1.04
		 & 0.0031
		 & 0.39
		 & 0.27
		 & 1.92                                                                                \\

		$\rho$
		 & 0.43
		 & 0.000024
		 & 0.0060
		 & --
		 & --                                                                                  \\
		\hline
	\end{tabular}
\end{table}
